% ============================================================
%  Georg Brandes, Levned, Bind I
%  Kjøbenhavn: Gyldendal, 1905.
%  Pp. 129–132 (om Hans Brøchner)
%  TRANSCRIPTION (Danish)
%  Source: author's copy of text.
% ============================================================
\documentclass[12pt,a4paper]{article}
\usepackage[utf8]{inputenc}
\usepackage[T1]{fontenc}
\usepackage[danish]{babel}
\usepackage{microtype}
\usepackage{setspace}
\usepackage[margin=1.2in]{geometry}
\usepackage{fancyhdr}
\usepackage{hyperref}

\hypersetup{
  pdftitle={Levned I — om Hans Brøchner},
  pdfauthor={Georg Brandes},
  hidelinks
}

\pagestyle{fancy}
\fancyhf{}
\rhead{Brandes, \emph{Levned} I (1905), pp.\ 129--132}
\cfoot{\thepage}

\setstretch{1.3}

% Original page-number marker
\newcommand{\pnum}[1]{\noindent\makebox[0pt][r]{\small\textit{[#1]}\quad}}

\title{\textbf{Levned}, Bind I\\[6pt]
  \large Uddrag om Hans Brøchner (pp.\ 129--132)}
\author{Georg Brandes}
\date{Kjøbenhavn: Gyldendal, 1905}

\begin{document}
\maketitle
\thispagestyle{fancy}

\bigskip

\noindent Længe havde jeg drevet mine ikke-juridiske Studier, som jeg
kunde bedst uden Vejledning af nogensomhelst Lærer. Men jeg havde følt
Savnet af en saadan. Og da en nyansat Docent ved Universitetet,
Professor H. Brøchner, tilbød Vejledning i Studiet af Filosofi til
hvem der vilde indfinde sig hos ham paa visse opgivne Tider, havde jeg
følt mig stærkt fristet til at benytte mig deraf. Jeg havde havt
Betænkeligheder derved, thi jeg vilde nødigt opgive det mindste af min
dyrebare Frihed; jeg nød min tilbagetrukne Tilstand, mit beskedne
Studielivs Hemmelighedsfuldhed; men paa den anden Side vilde jeg ikke
kunne give mig i Kast med Plato og Aristoteles uden Vink af en kyndig
Haand om hvad og hvorledes.

Jeg var da i stor Spænding. Jeg havde hørt Professor Brøchner tale om
Psykologi; men Foredraget var behandlet med en saa ængstelig Omhu, var
saa ensformigt og fremmedlydende, at det ikke kunde andet end skræmme.
Alene Professorens fornemme Skønhed tiltrak, især de store sorgfulde
Øjne med det dejlige Blik, der saa ud som havde de forsket og grædt.

\pnum{130}Nu besluttede jeg at vove mig op til Brøchner. Men jeg
havde ikke Mod til forud at tale derom til min Moder for ikke at
forskrække mig selv ved at nævne det; saa urolig var jeg ved Skridtet.
Da Dagen kom, jeg havde bestemt til Besøget --- det var den 2.\ September
1861 --- gik jeg flere Gange op og ned foran Huset, førend jeg kunde
bestemme mig til at gaa op ad Trappen; jeg søgte forud at beregne,
hvad Professoren vel vilde sige og hvad jeg da rettest skulde svare, i
det Uendelige.

Den høje, smukke Mand, der lignede en spansk Ridder, lukkede selv
Døren op og modtog venligt det unge Menneske, der snart skulde blive
hans nærmeste Elev. Han vidste til min Forundring, straks han hørte
mit Navn, fra hvilken Skole jeg var udgaaet og at jeg nylig var bleven
Student. Han fraraadede energisk en Gennemgaaen af Plato og Aristoteles
som altfor anstrengende --- sagde eller antydede, at jeg burde
forskaanes for den besværlige Vej, han selv havde tilbagelagt, og
skitserede en Studieplan for nyere Filosofi og Æstetik. Hans Holdning
indgav Mod og efterlod det Hovedindtryk, at han vilde spare Begynderen
enhver unyttig Anstrengelse. Jeg følte mig med mine unge Kræfter lige
fuldt, da jeg gik hjem, en Smule skuffet ved ikke at skulle tage fat
paa Filosofiens Historie forfra.

Besøget gentoges snart, og hurtigt udspandt der sig mellem Lærer og
Elev det inderligste Forhold, fra den ældres Side et faderligt
Velviljesforhold, som den yngre endnu ikke havde kendt, under hvilket
Læreren, undervisende og praktisk vejledende, uafbrudt havde
Opmærksomheden henvendt paa Alt, hvad der kunde fremme Lærlingens Vel
og sikre hans Fremtid; fra den yngres Side et Ærbødigheds- og
Hengivenhedsforhold, en dyb Taknemmelighed for den Omsorg, hvis
Genstand han var.

Jeg følte vel nu og da overfor denne Lærers store \pnum{131}Grundighed
og hans Dygtighed til at tumle de vanskeligste Tanker en pinlig
Mistillid til mine egne Evner og min egen Aandskraft i Sammenligning
med hans. Jeg følte mig heller ikke sjældent ilde berørt ved, at der
ligesom kom noget Skævt ind i vort Forhold, naar Brøchner tiltrods for
de Tvivl og Indvendinger, jeg fremsatte, altid tog som givet, at jeg
maatte dele hans panteistiske Grundanskuelse, uden Blik for, at jeg
endnu omtumledes i Tvivl og famlede efter et Fodfæste. Men det
fortrolige Forhold til den modnede Mand havde en Tiltrækning, som
Forholdet til de usikre eller ungdommeligt indskrænkede Kammerater
nødvendigvis maatte mangle: han havde et Livs Erfaring bag sig, han
saa fra oven nedad paa et ungt Menneskes Sympatier og Antipatier.

For mig stod f.\ Eks.\ dengang endnu Plougs \textit{Fædrelandet} som
Danmarks intelligente Organ, medens Billes \textit{Dagbladet} var mig
ledt, især paa Grund af den Overfladiskhed og den affærdigende Tone,
der betegnede Bladets literære Kritik. Brøchner, der med blandet
Velvilje og uden at gøre synderlig Forskel saa ned paa begge disse den
dannede Almenheds Blade, fandt i sin unge Lærlings Forbitrelse paa den
trivielle og dog nyttige Avis kun et Kendetegn paa hans
Sjælsbeskaffenhed. Men blot den Maade, paa hvilken Brøchner en Dag med
et Smil bemærkede: De læser nok af Princip ikke Dagbladet, fik mig til
i et Glimt at indse det Komiske i min Harme over en Hjerrilds Artikler.
Min Synskres var jo endnu snever nok til at Forholdene i Kjøbenhavn i
og for sig syntes mig at have Betydning. Selv fandt jeg min Synskres
omfattende. En Dag, da jeg gjorde mit indre Regnskab op, skrev jeg som
nittenaarig naivt: »Mine gode Egenskaber, de, der vil konstituere min
Personlighed, hvis jeg bliver til noget, er en vældig og glødende
Begejstring, en dygtig Myndighed i Sandhedens Tjeneste, en vid
\pnum{132}Synskres og en filosofisk dannet Tænkning. Det maa forsone
med min Mangel paa Humor og Lethed.«

I væsenlig Overensstemmelse med Hans Brøchner følte jeg mig først
flere Aar efter Bekendtskabets Stiftelse. Men da havde et Studium af
Ludwig Feuerbachs Skrifter henrevet mig; thi Feuerbach var den første
Tænker, hos hvem jeg fandt Gudsideens Opstaaen i det menneskelige Sind
fyldestgørende forklaret. Hos Feuerbach mødte der mig endelig en
Fremstilling uden Omsvøb og uden den tyske Filosofis sædvanlige tunge
Formler, en sejrrig Klarhed, som virkede højst velgørende paa mit eget
Tankesæt og gav mig Tryghed. Havde jeg i flere Aar følt mig betydeligt
tilhøjre for min Ven og Lærer, saa kom der nu en Tid, hvor jeg følte
mig adskilligt tilvenstre for ham med hans hemmelighedsfulde Leven i
det evige Aandsrige, af hvilket han følte sig som Led.

\end{document}
